% ============================================================
% 直击考场（杂志页码 72–75）
% ============================================================
\colchapter{直击考场}

\section*{在“人生三境”中拾级而上，方成大学问}
\auth{编写者：河北省唐山一中 吴国梁}

\subsection*{备考热题}

王国维曾说，古今之成大事业、大学问者，必经过三种境界：“昨夜西风凋碧树。独上高楼，望尽天涯路”，此第一境也。“衣带渐宽终不悔，为伊消得人憔悴”，此第二境也。“众里寻他千百度，蓦然回首，那人却在，灯火阑珊处”，此第三境也。

其实，做学问如此，社会生活和国家发展也是如此。无论个人的成长，还是成就一项伟大事业，甚至中国共产党的百年奋斗历程，无不是从这三种境界中一步步走过来的。

以上论述具有启示意义。请结合材料写一篇文章，体现你的感悟与思考。

要求：选准角度，确定立意，明确文体，自拟标题；不要套作，不得抄袭；不得泄露个人信息；不少于800字。

\subsection*{押题解读}

王国维的“人生三境界”直接体现了“成大事业、大学问者”需要经过的三层境界，分别阐释如下：

\textbf{第一境界为立志阶段}——树立高远目标，确定远大理想，不被眼前的苟且羁绊，不被安逸麻痹精神。

\textbf{第二境界为追求阶段}——坚忍不拔，坚毅执着，在萧索时潜得下心来，在低谷时耐得住寂寞。

\textbf{第三境界为收获阶段}——经过苦苦追求终有所悟，在某一瞬间豁然开朗，所有成就水到渠成。

王国维强调，“成大事业、大学问者”，这三个阶段缺一不可，且顺序不可乱，这就启示我们，高远的志向，不懈的奋斗，是成功的必要条件。此外，材料又进一步提示同学们“做学问如此，社会生活和国家发展也是如此”，这就要求把“大事业”“大学问”扩展到其他领域——可写人生的追求、时代的追求、历史的进步，可关注宏大叙事，也可着眼个体成长。总之，只要这些人与事的发展进程符合“立志——追求——收获”这三个阶段，就可成文。

在写作时，除了“三境界”外，还要注意“大”这一重点。“大”意味着眼界、胸怀、抱负、目标等的宏伟壮丽，其对立面就是偏狭、苟且、自私的人生态度，以及与之对应的投机取巧、自私自利、鼠目寸光、急功近利等不良品质，只有摒弃这些负面品质，才可以成其“大”。

\clearpage
\section*{做学问要向着“大师”努力}

向着“大师”努力，对应“三境界”中的第一境——立志。不论是做学问还是做人，都要立志高远，但成为“大师”既不是一朝一夕的事，也不是只属于某一个人生阶段的事，而应该是终身的事。在事业上取得极高造诣、极大成就，进而成为“大师”，这是一种“高远”，但还有另一种“高远”——学习不能只顾私人利益，求名应求万世名，不能局限于小家小团体，而要胸怀家国天下。这样的人相信“计利当计天下利，求名应求万世名”，能超脱小我小利，谋求大我大义，将个人与社会、时代结合起来，以大胸怀、大境界、大格局来求学、做事、做人，这样也能达到“大师”的境界，不愧于时代。

\begin{quoteBox}{热点素材}
1962年5月19日，曾任清华大学校长的梅贻琦先生溘然长逝。在其身后，人们这样评价他，“先生，粹然儒者、躬身身教，对国家之贡献独多且要。与并世诸君子比，华若不逮，而实则过之。”习近平总书记曾在与北京大学师生座谈时指出：“教师承担着最庄严、最神圣的使命。梅贻琦先生说：‘所谓大学者，非谓有大楼之谓也，有大师之谓也。’我体会，这样的大师，既是学问之师，又是品行之师。”梅贻琦曾在《大学一解》中把好的师生关系形容为“从游”，老师如大鱼，学生如小鱼，“大鱼前导，小鱼尾随，是从游也，从游既久，其濡染观摩之效，自不求而至，不为而成。”学生本来是小鱼，有了大鱼的指导，慢慢也会成长为大鱼。梅贻琦甘当人梯，以人格魅力引导学生，以学术造诣开启学生的智慧之门，是永远的“大师”。
\end{quoteBox}

\clearpage
\section*{现实生活中也存在境界层次不高的人}

人生意义，在于不断深化自己对社会、宇宙的认识。“不断深化”对应“三境界”中的第二境——追求。一个有鸿鹄之志的人，会从内心出发对外界进行探索，探索的边界越来越大，内心的认知也越来越丰富。然而，这个探索的过程恰恰是最熬炼人心境的，孤独、寂寞、茫然、彷徨、苦闷甚至自我怀疑等消极情绪都主要出现在这个阶段。曾国藩说：“凡事皆有极困极难之时，打得通的便是好汉。”这就是说，只有矢志不渝、坚持不懈探索世界并深化人生认识的人，才能到达最后的第三境。

\begin{quoteBox}{热点素材}
冯友兰认为，若我们能够了解人生，人生便有意义，反之，人生便无意义。而每个人对于人生的了解多不相同，因此，人的境界也存在分别。他根据人的认知差异，将人生的境界分为四类：自然境界，功利境界，道德境界，天地境界。

自然境界是最低级的。这一境界中的人，为人处世不是依照社会习惯，而是依照其本性，他们从未去了解做某种事情的意义。往好处说，这就是“天真烂漫”，往差处说便是“糊里糊涂”。这一类人，对于“生”“死”皆不了解，而且没有“我”的观念。

功利境界中的人对于人生的了解，比自然境界中的人进了一步，他们有“我”的观念，不论做什么事都是为了自己的利益。这类人大抵贪生怕死，就算为社会服务，行事的动机还是为小我换取更多利益。而道德境界则与功利境界完全相反，这一境界中的人，做事皆以服务社会为目的，他们能认识到，这世界上除了“我”以外，还有一个社会，一个整体。古今圣贤及英雄便已达到这一境界。

天地境界中的人，一切行动皆以服务宇宙为目的。这个“宇宙”指的不是天文学概念，而是指哲学意义上的“一切”。人的身躯虽小，但精神可充塞于天地之间，其行为不仅能服务于人类社会，更能对超出人类社会的整个宇宙有益，最终“与天地比寿，与日月同光”。唯大圣大贤能达到这个境界。
\end{quoteBox}

\subsection*{反面论证：现实中的境界层次}

这个角度适用于反面论证。现实生活中，有立志高远的人，也有蝇营狗苟的人；有坚定不移的人，也有半途而废的人；有达到崇高境界的人，当然也有“躺平”“摆烂”的人。境界层次低的人往往庸俗、世故、狡猾、贪婪，毫无道德底线，我们一方面要警惕自己堕落成反面教材，另一方面也要远离这样的人。

\begin{quoteBox}{热点素材}
“精致的利己主义者”一词来源于北京大学中文系教授钱理群的一段话：“我们的一些大学，包括北京大学，正在培养一些‘精致的利己主义者’，他们高智商，世俗，老道，善于表演，懂得配合，更善于利用体制达到自己的目的。这种人一旦掌握权力，比一般的贪官污吏危害更大。”精致的利己主义者，其一举一动都笼罩着利益的影子，或是把利益当作唯一驱动，或是吃透规则甚至“潜规则”，善于钻空子、找漏洞达到自身目的。鲁迅则认为这样的人像白蚁：“只顾自己一个的便利和舒服，凡有被他遇见的，都用作生活的材料，一路吃过去”“遗留下来的，却只是一条排泄的粪”……

我们要知道，不是一切都能用物质利益来衡量。白芳礼帮助贫困孩子实现梦想的爱意与怜惜，如何计算？杨善洲凝视曾经的荒山已是瓜果飘香，那份满足与充实，如何计算？张桂梅将大山女孩送往广阔世界时的感动与欣喜，又如何计算？帮助他人的快乐、承担责任的充实、坚守理想的笃定，才可能让我们拥有更丰富更完美的人生。
\end{quoteBox}

\clearpage
\section*{备考范本：攀登应须毕生力，风光贵在绝顶时}
\auth{文/河北省唐山一中高二（6）班 陈潇琦}

“昨夜西风凋碧树。独上高楼，望尽天涯路”“衣带渐宽终不悔，为伊消得人憔悴”“众里寻他千百度，蓦然回首，那人却在，灯火阑珊处”，这是王国维所说的“古今之成大事业、大学问者”所必经的三种境界。的确，无论是个人还是国家，正是因为走过了这三种境界，才能锤炼出厚重的生命，迎来绵长的太平岁月。于我而言，这三境界便是在探索目标中立志、在不断奋斗中磨炼耐心、在无数次回望中静悟的良性递进。

何为立志？就是俯瞰世间千万种选择后，以一种苏世独立，横而不流的志向与勇气，走上那条自己热爱的路，哪怕道阻且长。千百年来，无数仁人志士登高望远，阅尽天涯路后仍舍弃小我利益，走上满是坎坷艰辛的报国之路。古有文天祥“人生自古谁无死，留取丹心照汗青”的舍生取义，近有秋瑾“一腔热血勤珍重，洒去犹能化碧涛”的丹心赤诚。他们立下为国报效之志，不因金钱的诱惑和个人的安危灰心动摇，始终以自己的方式践行最初的诺言。立志要如山，方能坚定不移。

何为耐心？就是拨雪寻春、烧灯续昼的坚持。找到目标固然重要，但有耐心地专注于此，能让我们走得更远。屠呦呦一心抵抗疟疾，但真正使她发现青蒿素的，不只是这份决心，更是一天天重复实验、一次次翻阅医书典籍的耐心；黄旭华立志为国攻关，日复一日埋头于书案，终于用千万个数据“堆”出新中国第一艘核潜艇。这些成就，无疑是由奋斗与耐心共同造就的。我们要仰望目标，但更要低头专注于泥泞的路，一步一个脚印耐心攀登。唯有脚踏实地，砥砺前行，追梦人才能真正触碰到理想的曙光。

何为静悟？就是过尽千帆抵达彼岸后的思考与沉淀。我们曾跨越千山万水寻找答案，却在不经意间融会贯通，踏入灯火阑珊处。三毛说过：“学问之道，是人格的建立、生命的领悟、凡事广涵的体认——而不是做一架‘念书机器’。”读书需要我们慢慢参悟，生活也是如此，功夫到时自有答案，不必苛求。而达成所愿也不是就此万事大吉，我们还需要继续感悟世间百态，总结经验提升自我，向着人生更远处漫溯。

有志向，守初心，省自身，方能登上这三境界的高塔，赏会当凌绝顶之景色。

\begin{quoteBox}{名师点评}
作为一篇优秀的考场作文，本文标题巧妙化用著名数学家苏步青的名言“为学应须毕生力，攀登贵在少年时”，语言富有节奏感且内涵深刻。文章主体对“三境界”的递进关系进行了论证，观照而清晰地阐述了自己的理解框架，并简练概括核心词——“立志”“耐心”“静悟”，在论述过程中，作者综合运用典型人物素材与名言，在增强文章内涵的同时，也有效展示了自身丰厚的积累。
\end{quoteBox}
